\documentclass[10pt,a4paper]{amsart}
\usepackage[margin=19mm]{geometry}
\usepackage{amsmath,amssymb,mathtools}
\usepackage[scr=rsfs,cal=euler]{mathalfa}
\usepackage{xfrac}
\usepackage[unicode,hidelinks,bookmarks=false]{hyperref}
\newcommand{\ie}{i.e.\ }
\newcommand{\cf}{cf.\ }
\newcommand{\resp}{respectively\ }
\newcommand{\Subsection}{\subsection}
\providecommand{\nobreakdash}{\nobreak\hbox{-}}
\setlength{\emergencystretch}{2em}
\multlinegap0pt
\usepackage{enumerate}
\usepackage{xcolor}
\newtheorem{prop}{Proposition}
\title{Linear maps in $\mathcal{L}(\ell_{\MakeLowercase{p}},\mathcal{Y}) $ preserving parallel and TEA pairs}
\author{Arpita Mal}
\date{Source: arXiv:2606.03281v1 (CC BY 4.0)}
\begin{document}
\maketitle

\noindent\emph{Source and licence.} Linear maps in $\mathcal{L}(\ell_{\MakeLowercase{p}},\mathcal{Y}) $ preserving parallel and TEA pairs. Version arXiv:2606.03281v1 (CC BY 4.0), distributed by arXiv under the Creative Commons Attribution 4.0 International licence, http://creativecommons.org/licenses/by/4.0/, as recorded in the arXiv licence metadata for this exact version and shown on the abstract page https://arxiv.org/abs/2606.03281v1. Full licence notice: \texttt{licenses/arxiv-2606.03281-CC-BY-4.0.txt}.\par\smallskip

\noindent\textit{Editorial scope.} Counted unit: the article's prop prop-nec (source0.tex lines 815--828). The article's complete original proof of it is delivered with it (source0.tex lines 829--927, one \texttt{proof} environment of 99 lines). No mathematics is written, repaired or re-derived: the statement and the proof are the article's own lines, in the article's own notation, and no retained line is substituted or re-flowed.  No further source-local context is needed: every cross-reference of every retained line resolves inside the two delivered spans.  Nothing was omitted from any retained span, so the package carries no omission record.  No retained line contains a citation, so the wrapper carries no bibliography and no bibliography entry.  No retained line contains a figure environment, an image-inclusion command or drawing code, so no image is delivered and the package declares no asset dependency.  Editorial formatting only, in addition: an AMS \texttt{amsart} preamble that restates the article's own declarations and macros used by the retained lines, namely the environments \texttt{prop} on separate counters, together with the standard packages those lines need.  The retained environments print in this document as Proposition 1 in order of appearance, whereas the article prints the same items with the numbering of its own full document.  The complete native source of the article as distributed in the arXiv e-print for this exact version is bundled unchanged as \texttt{sources/201962/source0.tex}.
\begin{prop}\label{prop-nec}
	Suppose $(C_{11}\cap C_{22})\cup (C_{12}\cap C_{21})\neq \emptyset,$ and $T$ preserves parallel pairs. %, and  $C\neq \emptyset.$ 
	Then 
\begin{enumerate}
\item For each $\theta,$ $(\theta\in [0,2\pi)$ if $\mathbb{F}=\mathbb{C},$ and $\theta\in \{0,\pi\}$ if $\mathbb{F}=\mathbb{R})$ 
\begin{equation}\label{eq-06}
	\begin{split}
\max\{\cos(\theta+\phi_j):j\in C_{11}\cap C_{21}\}
\geq  \max\Big\{\alpha_s\cos(\theta+\phi_s):s\in C_{11}\cup C_{21}\Big\}.	
\end{split}
\end{equation}	
\item $|C_{11}\cap C_{21}|\geq 2.$
\end{enumerate}
\end{prop}

\begin{proof}
	
$(1)$  We first show that if $C_{12}\cap C_{21}\neq \emptyset,$ then $$\max\{\cos(\theta+\phi_j):j\in C_{11}\cap C_{21} \}
\geq \max\Big\{\alpha_i\cos(\theta+\phi_i): i\in C_{12}\cap C_{21}\Big\}.$$ 
Assume to the contrary that there exists $i\in C_{12}\cap C_{21}$ such that 
\[\max\{\cos(\theta+\phi_j):j\in C_{11}\cap C_{21} \}<\frac{|t_{i1}|}{k_1}\cos(\theta+\phi_i).\]	
Choose $n_0\in \mathbb{N}$ such that for all $j\in C_{11}\cap C_{21},$
\[\frac{2k_2n_0}{k_1}\Big\{\cos(\theta+\phi_j)-\frac{|t_{i1}|}{k_1}\cos(\theta+\phi_i)\Big\}<\frac{|t_{i1}|^2}{k_1^2}-1.\]
Consequently, for  $n_1\geq n_0,$
\begin{equation}\label{eq-01}
	\begin{aligned}
k_1^2+2k_1k_2n_1\cos(\theta+\phi_j)+k_2^2n_1^2&<|t_{i1}|^2+2	k_2n_1|t_{i1}|\cos(\theta+\phi_i)+k_2^2n_1^2\\
\Rightarrow|k_1+n_1k_2e^{i(\theta+\phi_j)}|&<||t_{i1}|+n_1k_2e^{i(\theta+\phi_i)}|\\
\Rightarrow|t_{j1}+n_1e^{i\theta}t_{j2}|&<|t_{i1}+n_1e^{i\theta}t_{i2}|.
\end{aligned}
\end{equation}
Choose 
$n_2\in \mathbb{N}$ sufficiently large so that $n_2>n_0,$ and for all $j\in C_{11}\cap C_{22},$
\begin{eqnarray*}
k_1^2-|t_{i1}|^2+n_2^2(|t_{j2}|^2-k_2^2)&<& 2n_2\Big(k_2	|t_{i1}|\cos(\theta+\phi_i)-k_1|t_{j2}|\cos(\theta+\phi_j)\Big).
\end{eqnarray*}
Note that such a choice of $n_2$ is possible, since $(|t_{j2}|^2-k_2^2)<0,$ for all $j\in C_{11}\cap C_{22}.$ Therefore, 
\begin{equation}\label{eq-02}
%	\begin{aligned}
		\begin{split}
	k_1^2+2n_2k_1|t_{j2}|\cos(\theta+\phi_j)+n_2^2|t_{j2}|^2&< |t_{i1}|^2+2n_2k_2	|t_{i1}|\cos(\theta+\phi_i)+n_2^2k_2^2\\
	\Rightarrow |k_1+n_2|t_{j2}|e^{i(\theta+\phi_j)}|&<||t_{i1}|+n_2k_2e^{i(\theta+\phi_i)}|\\
	\Rightarrow|t_{j1}+n_2e^{i\theta}t_{j2}|&<|t_{i1}+n_2e^{i\theta}t_{i2}|.
	\end{split}
%	\end{aligned}
\end{equation}
From \eqref{eq-01} and \eqref{eq-02}, we get, for all $j\in C_{11},$
\[|t_{j1}+n_2e^{i\theta}t_{j2}|<|t_{i1}+n_2e^{i\theta}t_{i2}|.\]
Therefore, $T(1,n_2e^{i\theta})$ does not attain its norm at the $j$-th coordinate, where $j\in C_{11}.$  So $\Big(T(1,n_2e^{i\theta}),T(1,0)\Big)$ is not a parallel pair in $\ell_\infty^m,$ whereas, $\Big((1,n_2e^{i\theta}),(1,0)\Big)$ is a parallel pair in $\ell_1^2.$ This contradiction proves that 
\[\max\{\cos(\theta+\phi_j):j\in C_{11}\cap C_{21}\}\\
\geq \max\Big\{\alpha_i\cos(\theta+\phi_i):i\in C_{12}\cap C_{21}\Big\}.\]
Next we show that  if $C_{11}\cap C_{22}\neq \emptyset,$ then $$\max\{\cos(\theta+\phi_j):j\in C_{11}\cap C_{21}\}
\geq \max\Big\{\alpha_l\cos(\theta+\phi_l): l\in C_{11}\cap C_{22}\Big\}.$$ Assume to the contrary that there exists $l\in C_{11}\cap C_{22}$ such that 
\[\max\{\cos(\theta+\phi_j):j\in C_{11}\cap C_{21}\}<\frac{|t_{l2}|}{k_2}\cos(\theta+\phi_l).\]	
Choose $m_0\in \mathbb{N}$ such that for all $j\in C_{11}\cap C_{21},$
\[\frac{2k_1m_0}{k_2}\Big\{\cos(\theta+\phi_j)-\frac{|t_{l2}|}{k_2}\cos(\theta+\phi_l)\Big\}<\frac{|t_{l2}|^2}{k_2^2}-1,\]
and for all $j\in C_{12}\cap C_{21},$
\begin{eqnarray*}
	m_0^2(|t_{j1}|^2-k_1^2)+(k_2^2-|t_{l2}|^2)+2m_0\Big(k_2	|t_{j1}|\cos(\theta+\phi_j)-k_1|t_{l2}|\cos(\theta+\phi_l)\Big)<0.
\end{eqnarray*}
Proceeding as previous, it is easy to check that  for all $j\in C_{21},$
\[|m_0t_{j1}+e^{i\theta}t_{j2}|<|m_0t_{l1}+e^{i\theta}t_{l2}|.\]
Therefore, $T(m_0,e^{i\theta})$ does not attain its norm at the $j$-th coordinate, where $j\in C_{21},$ and so $\Big(T(m_0,e^{i\theta}),T(0,1)\Big)$ is not a parallel pair in $\ell_\infty^m,$ whereas, $\Big((m_0,e^{i\theta}),(0,1)\Big)$ is a parallel pair in $\ell_1^2.$ This contradiction completes the proof.
%Choose $m_0\in \mathbb{N}$ such that for all $j\in C_{11}\cap C_{21},$
%\[\frac{2k_1m_0}{k_2}\Big\{\cos(\theta+\phi_j)-\frac{|t_{l2}|}{k_2}\cos(\theta+\phi_l)\Big\}<\frac{|t_{l2}|^2}{k_2^2}-1.\]
%Consequently, for  $m_1\geq m_0,$
%\begin{equation}\label{eq-03}
%	\begin{aligned}
%		m_1^2k_1^2+2k_1k_2m_1\cos(\theta+\phi_j)+k_2^2&<m_1^2k_1^2+2	k_1m_1|t_{l2}|\cos(\theta+\phi_l)+|t_{l2}|^2\\
%		\Rightarrow|m_1k_1+k_2e^{i(\theta+\phi_j)}|&<|m_1k_1+|t_{l2}|e^{i(\theta+\phi_l)}|\\
%		\Rightarrow|m_1t_{j1}+e^{i\theta}t_{j2}|&<|m_1t_{l1}+e^{i\theta}t_{l2}|.
%	\end{aligned}
%\end{equation}
%Choose 
%$m_2\in \mathbb{N}$ sufficiently large so that $m_2>m_0,$ and for all $j\in C_{12}\cap C_{21},$
%\begin{eqnarray*}
%m_2^2(|t_{j1}|^2-k_1^2)+(k_2^2-|t_{l2}|^2)+2m_2\Big(k_2	|t_{j1}|\cos(\theta+\phi_j)-k_1|t_{l2}|\cos(\theta+\phi_l)\Big)<0.
%\end{eqnarray*}
%Note that such a choice of $m_2$ is possible, since $(|t_{j1}|^2-k_1^2)<0,$ for all $j\in C_{12}\cap C_{21}.$ Therefore, 
%\begin{equation}\label{eq-04}
%	%	\begin{aligned}
%		\begin{split}
%			m_2^2|t_{j1}|^2+2m_2k_2|t_{j1}|\cos(\theta+\phi_j)+k_2^2&<m_2^2k_1^2 +2m_2k_1	|t_{l2}|\cos(\theta+\phi_l)+|t_{l2}|^2\\
%			\Rightarrow |m_2|t_{j1}|+k_2e^{i(\theta+\phi_j)}|&<|m_2k_1+|t_{l2}|e^{i(\theta+\phi_l)}|\\
%			\Rightarrow|m_2t_{j1}+e^{i\theta}t_{j2}|&<|m_2t_{l1}+e^{i\theta}t_{l2}|.
%		\end{split}
%		%	\end{aligned}
%\end{equation}
%From \eqref{eq-03} and \eqref{eq-04}, we get, for all $j\in C_{21},$
%\[|m_2t_{j1}+e^{i\theta}t_{j2}|<|m_2t_{l1}+e^{i\theta}t_{l2}|.\]
%Therefore, $T(m_2,e^{i\theta})$ does not attain its norm at the $j$-th coordinate, where $j\in C_{21},$ and so $\Big(T(m_2,e^{i\theta}),T(0,1)\Big)$ is not a parallel pair in $\ell_\infty^m,$ whereas, $\Big((m_2,e^{i\theta}),(0,1)\Big)$ is a parallel pair in $\ell_1^2.$ This contradiction proves that 
%\[\max\{\cos(\theta+\phi_j):j\in C_{11}\cap C_{21}\}\\
%\geq \max\Big\{\alpha_l\cos(\theta+\phi_l):l\in C_{11}\cap C_{22}\Big\}.\]

$(2)$ 
Suppose for a contradiction that  $C_{11}\cap C_{21}= \{i\}.$ \\
For $\mathbb{F}=\mathbb{R},$  choose $\theta\in \{0,\pi\}$ such that $\cos(\theta+\phi_i)=-1.$ For this choice of $\theta,$ \eqref{eq-06} is not satisfied, proving that   $|C_{11}\cap C_{21}|\geq 2.$\\
Suppose $\mathbb{F}=\mathbb{C}.$ Let $s\in(C_{11}\cap C_{22})\cup (C_{12}\cap C_{21}).$ Then at least two of the sets $A_s^+,A_s^-,A_s^0$ are empty. First assume $A_s^+=A_s^-=\emptyset,$ that is, $A_s^0=\{i\}.$ Clearly $\sin(\phi_i)\neq \alpha_s\sin(\phi_s),$ for otherwise, $\alpha_s=1$ contradicting that $s\in C.$ 
Choose $\theta\in(0,2\pi)$ such that $$\sin(\theta)\Big(\sin(\phi_i)-\alpha_s\sin(\phi_s)\Big)>0 .$$ Then  
\begin{eqnarray*}
	\cos(\theta)\Big(\cos(\phi_i)-\alpha_s \cos(\phi_s)\Big)=0&<&\sin(\theta) \Big(\sin(\phi_i)-\alpha_s \sin(\phi_s)\Big)\\
	\Rightarrow \cos(\theta-\phi_i)&<&\alpha_s\cos(\theta-\phi_s), %\\
%	&<& \max \{\alpha_s\cos(\theta-\phi_s):s\in C\},
\end{eqnarray*} 
which contradicts \eqref{eq-06}.  Now assume $A_s^+=A_s^0=\emptyset.$ Then $A_s^-=\{i\}.$ Choose $\theta\in (0,\pi)$ such that 
\begin{eqnarray*}
	\cos(\theta)&>& \sin(\theta) \frac{\sin(\phi_i)-\alpha_s \sin(\phi_s)}{\cos(\phi_i)-\alpha_s \cos(\phi_s)}\\
%	\Rightarrow  \cos(\theta)\Big(\cos(\phi_i)-\alpha_s \cos(\phi_s)\Big)&<&\sin(\theta) \Big(\sin(\phi_i)-\alpha_s \sin(\phi_s)\Big)\\
	\Rightarrow \cos(\theta-\phi_i)&<&\alpha_s\cos(\theta-\phi_s),%\\
%	&<& \max \{\alpha_s\cos(\theta-\phi_s):s\in C\},
\end{eqnarray*}
which again contradicts \eqref{eq-06}. Similarly $A_s^-=A_s^0=\emptyset$ produces a contradiction. Therefore, at most one of the sets $A_s^+,A_s^-,A_s^0$ is empty, and so $|C_{11}\cap C_{21}|\geq 2.$\\

\end{proof}

\end{document}
